\chapter{Kesimpulan dan Saran}

\section{Kesimpulan}
Berdasarkan serangkaian perancangan, implementasi mandiri dari awal (\textit{from scratch}), serta evaluasi empiris menggunakan protokol baku 5-\textit{fold cross-validation} pada dataset MovieLens 100K, kesimpulan penelitian ini dirumuskan untuk menjawab keempat pertanyaan penelitian:
\begin{enumerate}
    \item \textbf{Akurasi Prediksi Rating}: Ketiga algoritma sistem rekomendasi berbasis \textit{collaborative filtering} terbukti melampaui seluruh model \textit{baseline} secara signifikan. Model \textit{baseline} terkuat diraih oleh Rata-rata Film (MAE 0,8174; RMSE 1,0246; NDCG@10 0,8329). Metode \textit{User-Based Collaborative Filtering} (UBCF) menghasilkan tingkat kesalahan absolut terendah dengan MAE sebesar 0,7342 (peningkatan 10,18\% dibanding \textit{baseline}). Sementara itu, metode \textit{Matrix Factorization} (MF) mencatat kesalahan kuadratik terendah dengan RMSE sebesar 0,9249 (peningkatan 9,73\% dibanding \textit{baseline}), disusul oleh \textit{Item-Based Collaborative Filtering} (IBCF) dengan MAE 0,7383 dan RMSE 0,9459.
    
    \item \textbf{Pengaruh Parameter Ketetanggaan ($k$)}: Analisis sensitivitas pada UBCF dan IBCF membuktikan bahwa ukuran lingkungan ketetanggaan optimal berada pada rentang $k \in [30, 40]$. Nilai $k$ yang terlalu sempit ($k < 10$) menghasilkan varians estimasi yang tinggi akibat keterbatasan sampel tetangga, sedangkan nilai $k$ yang terlalu luas ($k > 50$) menurunkan akurasi karena mulai menyertakan pengguna atau item yang memiliki derajat kemiripan rendah (mengintroduksi derau). Selain itu, penggunaan \textit{Adjusted Cosine Similarity} pada IBCF terbukti secara konsisten mereduksi kesalahan (penurunan MAE $\approx 0,022$) dibandingkan \textit{Cosine} standar karena berhasil mengompensasi disparitas skala rating pengguna.
    
    \item \textbf{Kualitas Perankingan Daftar Rekomendasi (NDCG@10)}: Dalam skenario perankingan daftar Top-10 rekomendasi, \textit{Matrix Factorization} terbukti sebagai metode terbaik dengan skor NDCG@10 mencapai 0,8578 (peningkatan 2,99\% dibanding \textit{baseline} terkuat), mengungguli UBCF (0,8474) dan IBCF (0,8375). Berdasarkan uji signifikansi statistik berpasangan (\textit{paired t-test}), keunggulan perankingan Matrix Factorization bersifat sangat signifikan secara saintifik pada seluruh lima \textit{fold} ($p = 0,0102$ terhadap UBCF dan $p = 0,0008$ terhadap IBCF).
    
    \item \textbf{Karakteristik Trade-off dan Efisiensi Komputasi}: Terdapat kompromi yang jelas antara tahapan pelatihan dan tahapan inferensi. Algoritma \textit{memory-based} (UBCF dan IBCF) memiliki waktu pelatihan yang sangat singkat ($< 0,22$ detik), namun membutuhkan waktu inferensi yang lebih besar untuk menghitung bobot tetangga secara dinamis. Sebaliknya, \textit{Matrix Factorization} memerlukan pelatihan awal offline selama 29,26 detik, namun menghadirkan latensi inferensi online yang luar biasa cepat (hanya 0,0148 detik untuk 20.000 prediksi, atau kurang dari 1 mikrodetik per rekomendasi). Dengan demikian, \textit{Matrix Factorization} merupakan pendekatan paling unggul dan direkomendasikan untuk implementasi sistem rekomendasi skala industri karena menghadirkan kombinasi optimal antara akurasi kuadratik, kualitas perankingan, dan latensi inferensi \textit{real-time}.
\end{enumerate}

\section{Saran Pengembangan Lanjutan}
Beberapa arah pengembangan lanjutan yang disarankan untuk penelitian di masa mendatang meliputi:
\begin{enumerate}
    \item \textbf{Integrasi Umpan Balik Implisit (\textit{Implicit Feedback})}: Mengembangkan model faktorisasi matriks lebih lanjut seperti SVD++ atau \textit{Bayesian Personalized Ranking} (BPR) guna memanfaatkan data interaksi tidak langsung (seperti frekuensi klik, durasi menonton, atau riwayat pencarian) yang memiliki volume jauh lebih masif dibanding rating eksplisit.
    \item \textbf{Penyertaan Informasi Kontekstual (\textit{Context-Aware})}: Memperkaya model rekomendasi dengan faktor kontekstual temporal (hari, jam menonton), kondisi sosial penonton, dan jenis perangkat (\textit{mobile} atau TV) guna menghadirkan rekomendasi dinamis yang adaptif terhadap situasi pengguna saat mengakses platform.
    \item \textbf{Arsitektur Hibrida untuk Mengatasi Masalah \textit{Cold-Start}}: Membangun model sistem rekomendasi hibrida (\textit{Hybrid Recommender System}) yang memadukan representasi laten kolaboratif dengan pemodelan konten teks (deskripsi sinopsis, genre, aktor, sutradara) menggunakan teknik pemrosesan bahasa alami (\textit{Natural Language Processing}) sehingga sistem tetap mampu memberikan rekomendasi berkualitas tinggi terhadap film baru yang belum memiliki riwayat rating.
\end{enumerate}
